\chapter{Einleitung}\label{einleitung}

\section{Problemstellung und
Zielsetzung}\label{problemstellung-und-zielsetzung}

„Ein vernünftiger Live-Ticker ist fast unmöglich zu schreiben. Denn sein
Konzept ist die komplette Überforderung. Des Autors. Des Schreibens. Und
der Wirklichkeit.'' Mit diesen Worten beschreibt der Journalist
Constantin Seibt das Grundproblem der Liveticker-Berichterstattung ---
eine Einschätzung, die im modernen digitalen Sportjournalismus bis heute
nichts an Aktualität verloren hat (zitiert nach Beils 2023, S. 56). Der
Liveticker hat sich zur dominanten Darstellungsform der
Echtzeit-Berichterstattung im Profifußball entwickelt und steht dabei
exemplarisch für einen Zielkonflikt zwischen \textbf{Geschwindigkeit und
Qualität}, der mit rein manuellen Mitteln nicht auflösbar ist.

Die vorliegende Arbeit identifiziert drei Problemdimensionen, die den
Einsatz neuer technologischer Lösungen notwendig machen:

\begin{itemize}
\tightlist
\item
  Der \textbf{operative Zeitdruck} bei der manuellen Texterstellung:
  Redakteure müssen Spielereignisse in Echtzeit beobachten, einordnen
  und in stilistisch kohärente Kurztexte überführen --- typischerweise
  innerhalb von 30 bis 120 Sekunden pro Eintrag (vgl. Hauser 2008,
  S.~3--4; Feldbeobachtung im Projektzeitraum, bestätigt durch
  Experteninterview, Kap.~2.4).
\item
  Die Herausforderungen der \textbf{Internationalisierung} durch
  Mehrsprachigkeit: Profifußballvereine bedienen zunehmend
  internationale Fanbases; deutschsprachige Redaktionen stehen vor der
  Anforderung, Inhalte parallel auf Englisch oder weiteren Sprachen
  bereitzustellen, ohne ausreichend mehrsprachiges Redaktionspersonal
  vorzuhalten.
\item
  Der \textbf{strukturelle Wandel} der Vereine zu eigenständigen
  Medienproduzenten mit White-Label-Anforderungen: Vereine betreiben
  eigene digitale Kanäle und Medienmarken, die Redaktionswerkzeuge mit
  konfigurierbarem vereinsspezifischem Ton und CI-konformer Gestaltung
  erfordern.
\end{itemize}

Diese Dimensionen werden in Kapitel 2 wissenschaftlich fundiert und in
konkrete Systemanforderungen überführt.

Das übergeordnete Ziel der vorliegenden Arbeit ist daher die Konzeption
und Implementierung eines produktionsfähigen Redaktionssystems, das
diesen Konflikt durch den gezielten Einsatz von \textbf{Large Language
Models} auflöst. Das System wird dabei in zwei Ausprägungen realisiert:
einer generischen Instanz für beliebige Vereine sowie einer
vereinsspezifischen White-Label-Instanz am Beispiel \textbf{Eintracht
Frankfurt}. Beide Ausprägungen implementieren dieselben drei
Betriebsmodi; im \texttt{coop}-Modus verbleiben dabei die finale
Entscheidungshoheit und publizistische Verantwortung vollständig beim
Menschen, während der \texttt{auto}-Modus bewusst als Kontrollbedingung
ohne redaktionelles Korrektiv realisiert wird.

Um dieses Ziel zu operationalisieren, definiert und evaluiert die Arbeit
drei Betriebsmodi:

\begin{enumerate}
\def\labelenumi{\arabic{enumi}.}
\tightlist
\item
  \textbf{Status quo (\texttt{manual}):} Vollmanuelle Erstellung als
  Vergleichsbasis.
\item
  \textbf{Vollautonom (\texttt{auto}):} KI-basierte Texterstellung und
  Publikation ohne menschliches Korrektiv.
\item
  \textbf{Hybrid (\texttt{coop}):} Kombination aus KI-generierten
  Textvorschlägen mit redaktioneller Prüfung und Freigabe.
\end{enumerate}



\section{Forschungsfrage}\label{forschungsfrage}

Aus der in Abschnitt~1.1 entwickelten Zielsetzung ergibt sich folgende
Forschungsfrage, die der vorliegenden Arbeit zugrunde liegt:

\begin{quote}
Wie lässt sich ein hybrides KI-gestütztes Redaktionssystem für die
Liveticker-Erstellung im Profifußball konzipieren und prototypisch
implementieren --- und inwiefern erfüllt das System die definierten
journalistischen Qualitätsanforderungen hinsichtlich Korrektheit,
Tonalität und Verständlichkeit?
\end{quote}

Die Frage ist konstruktiv-evaluativ angelegt: Der erste Teil adressiert
die \textbf{Gestaltungs- und Implementierungsdimension} im Sinne des
Design Science Research (Hevner et al.~2004) --- Konzeption,
Architekturentscheid und prototypische Realisierung eines
produktionsfähigen Systems. Der zweite Teil adressiert die
\textbf{Qualitätsdimension} --- die empirische Prüfung, ob das
System journalistische Mindestanforderungen an Korrektheit, Tonalität
und Verständlichkeit erfüllt. Als ergänzenden Befund wird untersucht,
welches Potenzial zur Reduktion der Publikationslatenz sich aus den
gemessenen Systemlatenzen und der Systemarchitektur ableiten lässt;
da ein kontrollierter Feldvergleich nicht durchführbar war, ist dieser
Befund nicht als primäres Evaluationsziel zu verstehen. Die
Operationalisierung erfolgt in Kapitel 1.3 methodisch und in Kapitel 6
vollständig.



\section{Methodisches Vorgehen}\label{methodisches-vorgehen}

Die Arbeit folgt einem konstruktiven Forschungsansatz im Sinne des
\textbf{Design Science Research} (Hevner et al.~2004): Ein
softwarebasiertes Artefakt wird als Antwort auf ein identifiziertes
Praxisproblem entworfen, implementiert und evaluiert. Der
Erkenntnisgewinn entsteht dabei nicht nur aus der Evaluation, sondern
auch aus dem Entwurfs- und Implementierungsprozess selbst.

Das methodische Vorgehen gliedert sich in fünf Phasen:

\begin{enumerate}
\def\labelenumi{\arabic{enumi}.}
\tightlist
\item
  \textbf{Problemanalyse und Anforderungserhebung:} Literaturgestützte
  Vertiefung der Problemdimensionen und Ableitung funktionaler sowie
  nicht-funktionaler Anforderungen.
\item
  \textbf{Technologierecherche:} Aufarbeitung des Stands der Technik in
  den relevanten Bereichen und Bewertung hinsichtlich der Eignung für
  das identifizierte Problem.
\item
  \textbf{Systemkonzeption:} Entwurf einer dreischichtigen
  Systemarchitektur mit Datenmodell, LLM-Pipeline und Prompt-Design.
\item
  \textbf{Implementierung:} Umsetzung als produktionsfähige Webanwendung
  mit vollständiger API, browserbasiertem Frontend, automatisiertem
  ETL-System und Cloud-Deployment.
\item
  \textbf{Evaluation:} Technische Qualitätssicherung, Evaluation der
  KI-Textgenerierung sowie systematischer Anforderungsabgleich.
\end{enumerate}

Die Evaluationsmethodik operationalisiert die Forschungsfrage entlang
zweier Dimensionen: Die \textbf{zeitliche Dimension} wird über die
Time-to-Publish-Metrik (TTP) gemessen, die den Zeitraum zwischen
Spielereignis und Veröffentlichung des Ticker-Eintrags erfasst. Die
\textbf{qualitative Dimension} wird über eine manuelle Bewertung auf den
Skalen Korrektheit, Tonalität und Verständlichkeit operationalisiert,
ergänzt durch ein \textbf{strukturiertes Experteninterview} mit einem
professionellen Sportredakteur zur Validierung der Systemeignung im
redaktionellen Arbeitskontext.

\section{Abgrenzung}\label{abgrenzung}

Das System geht bewusst über den Umfang eines typischen akademischen
Projekts hinaus und zielt auf Produktionsfähigkeit: Eine umfassende
automatisierte Testsuite, ein Cloud-Deployment und die Integration
realer Datenquellen dokumentieren diesen Anspruch. Diese
Praxisorientierung ist durch den beruflichen Kontext des Autors
motiviert: Als Mitarbeiter der \textbf{Stackwork GmbH} im IT-Bereich von
Eintracht Frankfurt entstand das System in direkter Kooperation mit den
fachlichen Anforderungen einer professionellen Redaktion. Das Ergebnis
ist ein vollständiges, eigenständig lauffähiges Redaktionssystem, das
als Cloud-Service alle Spiele der konfigurierten Wettbewerbe
verarbeitet. Zusätzlich werden publizierte Inhalte über Export-Workflows
an die bestehende \textbf{Stackwork Demo App} übertragen, wodurch das
Projekt zwei unabhängige Zielsysteme über nahezu identische Workflows
bedient. Die technischen Details der Architektur und Implementierung
werden in Kapitel 4 und 5 beschrieben.

Dennoch bestehen Einschränkungen hinsichtlich der Authentifizierung
(kein vollständiges Rollenmodell implementiert), der
Evaluationsunabhängigkeit (der Autor ist gleichzeitig Entwickler und
Bewerter) sowie der thematischen Reichweite (ausschließlich
Fußball-Liveticker), die in Kapitel~7 kritisch reflektiert werden.



\section{Aufbau der Arbeit}\label{aufbau-der-arbeit}

Die vorliegende Arbeit gliedert sich in acht Kapitel:

\begin{itemize}
\tightlist
\item
  \textbf{Kapitel 1 -- Einleitung} \emph{(vorliegendes Kapitel)}:
  Problemstellung, Forschungsfrage, methodisches Vorgehen und Aufbau der
  Arbeit.
\item
  \textbf{Kapitel 2 -- Motivation und Anforderungen:} Vertiefung der
  drei Problemdimensionen (redaktioneller Aufwand, Mehrsprachigkeit,
  White-Label-Bedarf), Ableitung der Systemanforderungen am Beispiel
  Eintracht Frankfurt, Experteninterview-Leitfaden und formale
  Anforderungsdefinition.
\item
  \textbf{Kapitel 3 -- Stand der Technik:} Large Language Models, Prompt
  Engineering, Natural Language Generation im Sport sowie
  Echtzeit-Technologien und ETL-Prozesse.
\item
  \textbf{Kapitel 4 -- Systemkonzeption:} Dreischichtige Architektur,
  Datenmodell, LLM-Pipeline und Prompt-Design.
\item
  \textbf{Kapitel 5 -- Implementierung:} Umsetzung als produktionsfähige
  Webanwendung mit über 70 API-Endpunkten, TypeScript-Frontend und
  n8n-ETL-System; ergänzt durch vollständige TypeScript-Typisierung,
  automatisierte Testsuite und Cloud-Deployment.
\item
  \textbf{Kapitel 6 -- Evaluation:} Technische Qualitätssicherung
  (Tests, Coverage, Typsicherheit), Evaluation der KI-Textgenerierung
  und systematischer Anforderungsabgleich.
\item
  \textbf{Kapitel 7 -- Diskussion:} Einordnung in den Stand der Technik,
  kritische Reflexion der Limitationen und Implikationen für den
  Sportjournalismus.
\item
  \textbf{Kapitel 8 -- Fazit und Ausblick:} Zusammenfassung der
  Erkenntnisse, Beantwortung der Forschungsfrage und Ausblick.
\end{itemize}
